\chapter{Computability and Complexity}\label{computability-and-complexity}
\chapterstatus{planned}

What can be computed at all, and what can be computed efficiently? This chapter makes the notion of an
algorithm precise in several equivalent ways, proves that some problems have no algorithmic solution, and
introduces the complexity classes that measure the cost of computation.

\noindent\textbf{Prerequisites.} Chapter~\ref{logic} (Mathematical Logic), Chapter~\ref{sets} (Set Theory) and Chapter~\ref{number-systems} (Number Systems).

\noindent\textbf{Planned contents.}
\begin{enumerate}
\item Primitive recursive and partial recursive functions; the Ackermann function is not primitive recursive.
\item Turing machines and register machines; the Church--Turing thesis; equivalence of the models; universal machines and the $s$-$m$-$n$ theorem.
\item Recursively enumerable sets; the halting problem; Rice's theorem; many-one and Turing reducibility; Kleene's recursion theorem.
\item Undecidable problems from mathematics: Hilbert's tenth problem (statement of the Matiyasevich--Robinson--Davis--Putnam theorem, with the proof of the exponential case), the word problem for groups, Post correspondence.
\item Connection with Chapter~\ref{logic}: representability of recursive functions in arithmetic and a second proof of the incompleteness theorems.
\item Time and space complexity; the classes $\mathrm{P}$, $\mathrm{NP}$, $\mathrm{PSPACE}$, $\mathrm{EXP}$; the hierarchy theorems; Savitch's theorem.
\item $\mathrm{NP}$-completeness: the Cook--Levin theorem, reductions to 3-SAT, clique, Hamilton cycle, subset sum; the $\mathrm{P}$ versus $\mathrm{NP}$ problem.
\item Randomised computation ($\mathrm{BPP}$), interactive proofs and an overview of $\mathrm{IP} = \mathrm{PSPACE}$; circuit complexity.
\item Exercises with solutions.
\end{enumerate}
